% GENERATED FILE. Edit canonical lessons, then run npm run generate:books.
% canonical-source-hash: fnv1a64:e1fe57d9fd144b98
% canonical-lessons: UR-C387-quli, UR-C387-sanduq, UR-C387-ladna, UR-C387-dhhona, UR-C387-bithana

\chapter{Porters and Loads}
\label{ch:ur-porters-and-loads}
\hlchaptermodality{387}

\begin{chapteropening}
\textbf{By the end of this chapter:} \emph{I can talk about porters, trunks and the loads people carry.}

Complete the last lesson of chapter 387: I can talk about porters, trunks and the loads people carry.
\end{chapteropening}

\section[\texorpdfstring{\ur{قلی} (qulī)}{قلی (qulī)}]{\ur{قلی} (qulī) — a porter}

\label{lesson:UR-C387-quli}



\begin{quote}
Before the new one: say the Urdu for a whistle, then the Urdu for to sound.
\end{quote}

\subsection*{You'll want to know: \ur{قلی}}
\textbf{\ur{قلی}} — \emph{qulī} — \textquotedblleft{}a porter\textquotedblright{}.

\begin{quote}
\ur{اسٹیشن} \ur{پر} \ur{قلی} \ur{لال} \ur{قمیض} \ur{پہنتے} \ur{ہیں۔}
\end{quote}

\begin{quote}
\emph{isṭeshan par qulī lāl qamīz pahante haiṅ.}
\end{quote}

\begin{quote}
At the station the porters wear red shirts.
\end{quote}

Masculine.

\begin{cousinweb}
\textbf{\ur{قلی}} (\emph{qulī}) is often traced to Tamil \emph{kūli}, \textquotedblleft{}wages, hire\textquotedblright{} (DEDR 1905): the porter is the man paid by the load. Other sources have been proposed, so the link is likely rather than certain.
\end{cousinweb}

\begin{grammarlens}[title={What you've built}]
One more word for getting around.
\end{grammarlens}

\subsection*{Your turn}
\begin{itemize}
\raggedright
  \item \emph{Say it:} \emph{qulī}
  \item \emph{Say it:} \emph{qulī}, once more
  \item \emph{Say it:} \emph{isṭeshan par qulī lāl qamīz pahante haiṅ.} — At the station the porters wear red shirts.
  \item \emph{Recall:} say the Urdu for a whistle, then the Urdu for to sound, then say \emph{qulī} again
\end{itemize}

\begin{tcolorbox}[breakable,colback=teal!4,colframe=teal!35!black,title={Before you move on}]
What does \ur{قلی} mean? (\textquotedblleft{}A porter\textquotedblright{}.) Say it once more.
\end{tcolorbox}
\section[\texorpdfstring{\ur{صندوق} (sandūq)}{صندوق (sandūq)}]{\ur{صندوق} (sandūq) — a trunk; a chest}

\label{lesson:UR-C387-sanduq}



\begin{quote}
Before the new one: say the Urdu for to sound, then the Urdu for a porter.
\end{quote}

\subsection*{You'll want to know: \ur{صندوق}}
\textbf{\ur{صندوق}} — \emph{sandūq} — \textquotedblleft{}a trunk; a chest\textquotedblright{}.

\begin{quote}
\ur{قلی} \ur{نے} \ur{میرا} \ur{صندوق} \ur{اٹھایا۔}
\end{quote}

\begin{quote}
\emph{qulī ne merā sandūq uṭhāyā.}
\end{quote}

\begin{quote}
The porter picked up my trunk.
\end{quote}

Masculine.

\begin{cousinweb}
\textbf{\ur{صندوق}} (\emph{sandūq}) is usually traced, through Persian, to Arabic \emph{ṣundūq}, \textquotedblleft{}a chest, a box\textquotedblright{}.
\end{cousinweb}

\begin{grammarlens}[title={What you've built}]
One more word for getting around.
\end{grammarlens}

\subsection*{Your turn}
\begin{itemize}
\raggedright
  \item \emph{Say it:} \emph{sandūq}
  \item \emph{Say it:} \emph{sandūq}, once more
  \item \emph{Say it:} \emph{qulī ne merā sandūq uṭhāyā.} — The porter picked up my trunk.
  \item \emph{Recall:} say the Urdu for to sound, then the Urdu for a porter, then say \emph{sandūq} again
\end{itemize}

\begin{tcolorbox}[breakable,colback=teal!4,colframe=teal!35!black,title={Before you move on}]
What does \ur{صندوق} mean? (\textquotedblleft{}A trunk; a chest\textquotedblright{}.) Say it once more.
\end{tcolorbox}
\section[\texorpdfstring{\ur{لادنا} (lādnā)}{لادنا (lādnā)}]{\ur{لادنا} (lādnā) — to load}

\label{lesson:UR-C387-ladna}



\begin{quote}
Before the new one: say the Urdu for a porter, then the Urdu for a trunk.
\end{quote}

\subsection*{You'll want to know: \ur{لادنا}}
\textbf{\ur{لادنا}} — \emph{lādnā} — \textquotedblleft{}to load\textquotedblright{}.

\begin{quote}
\ur{مزدور} \ur{ٹرک} \ur{پر} \ur{سامان} \ur{لادتے} \ur{ہیں۔}
\end{quote}

\begin{quote}
\emph{mazdūr ṭrak par sāmān lādte haiṅ.}
\end{quote}

\begin{quote}
The labourers load the goods onto the lorry.
\end{quote}

Past, with \ur{نے} (he or she …): \ur{اس} \ur{نے} \ur{گاڑی} \ur{پر} \ur{سامان} \ur{لادا}, \emph{us ne gāṛī par sāmān lādā}.

\begin{cousinweb}
\textbf{\ur{لادنا}} (\emph{lādnā}) is inherited from Sanskrit \emph{lardayati}, \textquotedblleft{}loads\textquotedblright{} (CDIAL 10966).
\end{cousinweb}

\begin{grammarlens}[title={What you've built}]
One more word for getting around.
\end{grammarlens}

\subsection*{Your turn}
\begin{itemize}
\raggedright
  \item \emph{Say it:} \emph{lādnā}
  \item \emph{Say it:} \emph{lādnā}, once more
  \item \emph{Say it:} \emph{mazdūr ṭrak par sāmān lādte haiṅ.} — The labourers load the goods onto the lorry.
  \item \emph{Recall:} say the Urdu for a porter, then the Urdu for a trunk, then say \emph{lādnā} again
\end{itemize}

\begin{tcolorbox}[breakable,colback=teal!4,colframe=teal!35!black,title={Before you move on}]
What does \ur{لادنا} mean? (\textquotedblleft{}To load\textquotedblright{}.) Say it once more.
\end{tcolorbox}
\section[\texorpdfstring{\ur{ڈھونا} (ḍhonā)}{ڈھونا (ḍhonā)}]{\ur{ڈھونا} (ḍhonā) — to haul; to carry (a heavy load)}

\label{lesson:UR-C387-dhhona}



\begin{quote}
Before the new one: say the Urdu for a trunk, then the Urdu for to load.
\end{quote}

\subsection*{You'll want to know: \ur{ڈھونا}}
\textbf{\ur{ڈھونا}} — \emph{ḍhonā} — \textquotedblleft{}to haul; to carry (a heavy load)\textquotedblright{}.

\begin{quote}
\ur{قلی} \ur{دن} \ur{بھر} \ur{بھاری} \ur{صندوق} \ur{ڈھوتے} \ur{ہیں۔}
\end{quote}

\begin{quote}
\emph{qulī din bhar bhārī sandūq ḍhote haiṅ.}
\end{quote}

\begin{quote}
Porters haul heavy trunks all day long.
\end{quote}

Past, with \ur{نے} (he or she …): \ur{اس} \ur{نے} \ur{سامان} \ur{ڈھویا}, \emph{us ne sāmān ḍhoyā}.

\begin{cousinweb}
\textbf{\ur{ڈھونا}} (\emph{ḍhonā}) is inherited: CDIAL lists the Hindi cousin \emph{ḍhonā}, \textquotedblleft{}to convey\textquotedblright{}, under Sanskrit \emph{ḍhaukayati}, \textquotedblleft{}brings near\textquotedblright{} (CDIAL 5610).
\end{cousinweb}

\begin{grammarlens}[title={What you've built}]
One more word for getting around.
\end{grammarlens}

\subsection*{Your turn}
\begin{itemize}
\raggedright
  \item \emph{Say it:} \emph{ḍhonā}
  \item \emph{Say it:} \emph{ḍhonā}, once more
  \item \emph{Say it:} \emph{qulī din bhar bhārī sandūq ḍhote haiṅ.} — Porters haul heavy trunks all day long.
  \item \emph{Recall:} say the Urdu for a trunk, then the Urdu for to load, then say \emph{ḍhonā} again
\end{itemize}

\begin{tcolorbox}[breakable,colback=teal!4,colframe=teal!35!black,title={Before you move on}]
What does \ur{ڈھونا} mean? (\textquotedblleft{}To haul; to carry (a heavy load)\textquotedblright{}.) Say it once more.
\end{tcolorbox}
\section[\texorpdfstring{\ur{بٹھانا} (biṭhānā)}{بٹھانا (biṭhānā)}]{\ur{بٹھانا} (biṭhānā) — to seat (someone)}

\label{lesson:UR-C387-bithana}



\begin{quote}
Before the new one: say the Urdu for to load, then the Urdu for to haul.
\end{quote}

\subsection*{You'll want to know: \ur{بٹھانا}}
\textbf{\ur{بٹھانا}} — \emph{biṭhānā} — \textquotedblleft{}to seat (someone)\textquotedblright{}.

\begin{quote}
\ur{قلی} \ur{نے} \ur{سارا} \ur{کو} \ur{کھڑکی} \ur{کے} \ur{پاس} \ur{بٹھایا۔}
\end{quote}

\begin{quote}
\emph{qulī ne sārā ko khiṛkī ke pās biṭhāyā.}
\end{quote}

\begin{quote}
The porter gave Sara a seat by the window.
\end{quote}

Past, with \ur{نے} (he or she …): \ur{اس} \ur{نے} \ur{بچے} \ur{کو} \ur{بٹھایا}, \emph{us ne bachche ko biṭhāyā}.

\begin{cousinweb}
\textbf{\ur{بٹھانا}} (\emph{biṭhānā}) is built on \textbf{\ur{بیٹھنا}}, \textquotedblleft{}to sit\textquotedblright{}, whose Hindi cousin \emph{baiṭhnā} CDIAL lists under Sanskrit \emph{upaviśati}, \textquotedblleft{}sits down\textquotedblright{} (CDIAL 2245): this is the form for making someone sit.
\end{cousinweb}

\begin{grammarlens}[title={What you've built}]
One more word for getting around.
\end{grammarlens}

\subsection*{Your turn}
\begin{itemize}
\raggedright
  \item \emph{Say it:} \emph{biṭhānā}
  \item \emph{Say it:} \emph{biṭhānā}, once more
  \item \emph{Say it:} \emph{qulī ne sārā ko khiṛkī ke pās biṭhāyā.} — The porter gave Sara a seat by the window.
  \item \emph{Recall:} say the Urdu for to load, then the Urdu for to haul, then say \emph{biṭhānā} again
\end{itemize}

\begin{tcolorbox}[breakable,colback=teal!4,colframe=teal!35!black,title={Before you move on}]
What does \ur{بٹھانا} mean? (\textquotedblleft{}To seat (someone)\textquotedblright{}.) Say it once more.
\end{tcolorbox}
